\section{Related Work}
\label{sec:related}

Transformer-based LLMs have demonstrated significant potential in numerous human-level language modeling tasks by continuously increasing their sizes~\cite{vaswani2017attention, rae2021scaling, smith2022using, du2022glam, chowdhery2022palm}.
%%%
As GPT-3 becomes the first model to surpass 100B parameters~\cite{brown2020language}, multiple LLMs ($>$100B) have been released, including OPT-175B~\cite{zhang2022opt}, Bloom-176B~\cite{scao2022bloom}, and PaLM~\cite{chowdhery2022palm}.
%%%
%Due to the auto-regressive nature of LLMs and their humongous model size, inference latency can be non-negligible in practice. 
%%%
Recent work has proposed a variety of approaches to accelerating generative LLM inference, which can be categorized into two classes.

\paragraph{Lossless acceleration.}
Prior work has explored the idea of using an LLM as a verifier instead of a decoder to boost inference.
%%%
For example, Yang et al.~\cite{yang2023inference} introduced {\em inference with reference}, which leverages the overlap between an LLM's output and the references obtained by retrieving documents, and checks each reference's appropriateness by examining the decoding results of the LLM.
Motivated by the idea of speculative execution in processor optimizations~\cite{burton1985speculative, hennessy2011computer}, recent work proposed {\em speculative decoding}, which uses a small language model to produce a sequence of tokens and examines the correctness of these tokens using an LLM~\cite{leviathan2022fast, xiaspeculative, stern2018blockwise, chen2023accelerating, gante2023assisted}.
%%%
There are two key differences between \Sys and these prior works.
First, instead of only considering a single sequence of tokens, \Sys generates and verifies a token tree, whose nodes each represent a unique token sequence. \Sys performs tree attention to compute the attention output of these token sequences in parallel and uses a novel tree-based decoding algorithm to reuse intermediate results shared across these sequences.
%%%
Second, prior attempts generally consider a single small language model for speculation, which cannot align well with an LLM due to the model capacity gap between them.
%%%
\rev{\Sys introduces two novel speculation methods, including 1) expanding from a single SSM and 2) merging from multiple fine-tuned SSMs, and the generated token tree largely increases the coverage of the LLM's output.}
% \Sys introduces collective boost-tuning to adapt different SSMs to align with an LLM under different scenarios, which largely increases the coverage of the speculated token trees produced by \Sys.
%Third, an additional challenge \Sys has to address is deciding the speculative configuration for a given request. \Sys leverages an important observation that the tokens generated by an LLM involve diverse difficulties to speculate, and uses a learning-based speculator to decide which SSMs to use and their speculative configurations.

Prior work has also introduced a variety of techniques to optimize ML computations on modern hardware platforms. 
%%%
For example, TVM~\cite{chen2018tvm} and Ansor~\cite{zheng2020ansor} automatically generate kernels for a given tensor program. TASO~\cite{jia2019taso} and PET~\cite{wang2021pet} automatically discover graph-level transformations to optimize the computation graph of a DNN.
\Sys's techniques are orthogonal and can be combined with these systems to accelerate generative LLM computation, which we believe is a promising avenue for future work.

% \ZJ{ZZ's original version:}
% In order to preserve output equivalency while reducing inference latency, \cite{leviathan2022fast, yang2023inference, xiaspeculative, stern2018blockwise, chen2023accelerating} bring the idea of speculative execution in processor optimizations~\cite{burton1985speculative, hennessy2011computer} to LLM inference optimization. 
% %%%
% Instead of solely relying on the LLM for inference, this line of work uses cheaper and faster proxies to do speculations followed by a verification procedure using the slower LLM to preserve equivalency.
% %%%
% Nevertheless, the quality of the speculation proxies as well as the system optimization opportunities for the verification process have been underexplored, which are the key focuses of this work.  


\paragraph{Lossy acceleration.}
BiLD~\cite{biglittletransformer} is a speculative decoding framework that uses a single SSM to accelerate LLM decoding. Unlike the systems mentioned above, the acceleration is lossy: speed comes at the cost of a possible degradation in the generated tokens. 
Another line of research leverages model compression to reduce LLM inference latency while compromising the predictive performance of the LLM.
%%%
% To this end, recent works have introduced a variety of techniques to reduce the end-to-end latency of serving LLMs. 
%%%
For example, prior work proposed to leverage weight/activation quantization of LLMs to reduce the memory and computation requirements of serving these LLMs~\cite{xiao2022smoothquant, frantar2023gptq, park2022nuqmm, yao2022zeroquant,dettmers2022gpt3}.
%%%
Recent work further explores a variety of structured pruning techniques for accelerating Transformer-based architectures~\cite{frantar2023sparsegpt, wang2021spatten, hubara2021accelerated}.
%%%
%Regardless, few of them have been used in practice due to unpredictable performance degradation. 
%%%
A key difference between \Sys and these prior works is that \Sys does not directly reduce the computation requirement for performing LLM inference, but instead reorganizing LLM inference computation in a more parallelizable way, which reduces memory accesses and inference latency at the cost of manageable memory and computation overheads.

\paragraph{Tree-structured attention.}
%Prior work by Nguyen et al.~\cite{Nguyen2020} introduced {\em tree-structured attention} , a technique to let a Transformer model capture the hierarchical composition of the input text by running the model on the text's parse tree. The key idea is to use a one-on-one mapping called {\em hierarchical accumulation}, to encode the parse tree in a form that can be processed by the attention, and then decode it back into tree form.
%%%
%There are two key difference between the tree-structured attention and \sys's tree-based decoding mechanism. First, \sys employs a tree structure to combine multiple candidate sequences together so as to condense identical prefixes; whereas Nguyen et al. represent a single input sequence using its parse tree. \sys does not attempt to incorporate into the LLM model additional information about each input sequence's structure, but rather, it accelerates LLM inference by verifying multiple decoded sequences (each generated by a SSM) in parallel. Second, unlike in the work by Nguyen et al., \sys's attention does not output another tree object, but rather, a single sequence of tokens.
Nguyen et al.~\cite{Nguyen2020} introduced {\em tree-structured attention}, a technique that lets a Transformer model capture the hierarchical composition of input text by running the model on the text's parse tree. To process with attention, it uses a one-on-one mapping to encode and decode the tree.
% , so that the attention can process it. 
There are two key differences from \sys's tree-based decoding. First, \sys uses a tree to combine candidate sequences to condense prefixes, whereas Nguyen et al. represent a single sequence with its parse tree. \sys does not incorporate parse tree into the LLM, but accelerates inference by verifying decoded sequences in parallel. Second, \sys's attention outputs a token sequence, not a tree.

\paragraph{Multi-sample decoding techniques}
% Like tree-based speculative inference, multi-sample decoding techniques such as \textit{beam search}, \textit{top-k sampling}, and \textit{top-p sampling} consider multiple candidate token sequences or next tokens at each step, rather than just the single highest probability option. They can also prune or ignore low probability sequences or tokens. There are two key differences from \sys's tree-based decoding. First, tree-based decoding in \sys is used to speculatively predict and verify multiple candidate token sequences in parallel against a large language model (LLM). The goal is to reduce latency by decreasing the number of decoding iterations needed. Beam search and top-k/top-p sampling are used during the decoding process to generate high-probability token sequences. They do not decrease the number of decoding steps.  Second, tree-based decoding leverages small speculative models (SSMs) to predict candidate sequences verified in parallel against the LLM. Beam search and top-k/top-p sampling are decoding strategies applied directly to the probability distribution over next tokens provided by the LLM itself.
% Overall, beam search, top-k sampling, and top-p sampling are decoding optimizations that are orthogonal to the techniques proposed in this paper. \sys supports using these techniques, which can be combined with tree-based speculative decoding.
Like tree-based speculative inference, \textit{beam search}, \textit{top-k sampling}, and \textit{top-p sampling} consider multiple candidate token sequences at each step and can prune low-probability options. However, tree-based decoding in \sys speculatively predicts and verifies multiple candidates in parallel against an LLM to reduce decoding iterations and latency, leveraging small speculative models (SSMs). In contrast, beam search and top-k/top-p sampling are decoding strategies applied directly to the LLM's output probabilities to generate high-probability sequences without reducing decoding steps. \sys supports beam search, top-k sampling, and top-p sampling. These techniques are orthogonal decoding optimizations and can be combined with tree-based speculative decoding.
 
